\documentclass[a4paper, 12pt]{article}
\usepackage[T2A]{fontenc}
\usepackage[utf8]{inputenc}
\usepackage[russian]{babel}
\usepackage{amsmath, amssymb}
\usepackage{physics}
\usepackage{bm}
\usepackage{braket}
\usepackage{hyperref}
\usepackage{graphicx}
\usepackage{xcolor}
\usepackage{subcaption}
\usepackage[verbose=true, a4paper, left=20mm, right=20mm, top=20mm, bottom=20mm]{geometry}
\usepackage{tikz}
\usetikzlibrary{decorations.pathmorphing}
\usetikzlibrary{decorations.markings}
\usepackage{pgfplots}
\pgfplotsset{compat=1.16}
\usepgfplotslibrary{groupplots}
\usepackage[backend=biber, style=phys, biblabel=brackets, chaptertitle=false,
pageranges=false]{biblatex}
\addbibresource{qm1_lecture01.bib}

\title{Конспект лекций по курсу <<Квантовая Механика 1>>  \\ \large Лекция 1}
\author{Лектор: Грибанов Сергей Сергеевич}
\date{1 сентября}

\begin{document}

\maketitle

\tableofcontents
\section{Введение}

Квантовая механика изучает поведение микроскопических объектов --- элементарных
частиц, атомов, молекул и других квантовых систем. Классическая физика,
включающая механику Ньютона и электродинамику Максвелла, оказалась неспособной
объяснить целый ряд экспериментальных фактов, относящихся к явлениям атомного
масштаба. К концу XIX века накопились экспериментальные данные, которые не укладывались в рамки
классических представлений: спектр излучения абсолютно чёрного тела,
фотоэффект, устройство атомных спектров, устойчивость атомов. Попытки
описать эти явления на основе классических законов приводили к противоречиям
с опытом.

Поведение объектов микромира носит принципиально вероятностный характер. Это
связано с тем, что электроны, протоны, нейтроны и другие частицы обладают не
только корпускулярными, но и волновыми свойствами. Для них, как и для
электромагнитного излучения, справедливы соотношения неопределённостей.
Поэтому в квантовой механике, в отличие от классической, понятие траектории
неприменимо. Вероятностная природа законов движения микроскопических объектов
не сводится к незнанию точных значений скрытых параметров, как в
статистической физике, а является их собственным свойством. Кроме того,
некоторые физические величины, характеризующие квантовые объекты, могут
изменяться не непрерывно, а дискретно --- определёнными порциями (квантами).

Квантовая механика объясняет строение атомов и молекул, природу химической
связи, свойства твёрдых тел, явления радиоактивного распада и многие другие
процессы, происходящие в микромире. Она лежит в основе современной атомной и
ядерной физики, физики элементарных частиц и физики конденсированного
состояния, а также находит применение в астрофизике и других областях.

\section{Предпосылки создания квантовой механики}

Изложение курса начинается с рассмотрения экспериментальных предпосылок,
которые привели к созданию квантовой механики. К ним, например, относятся
спектр излучения абсолютно чёрного тела, фотоэффект, эффект Комптона, опыт
Резерфорда и атомные спектры. Эти явления послужили основой для формирования
новых представлений о движении и состоянии микрообъектов.

\subsection{Излучение абсолютно чёрного тела}

Абсолютно чёрное тело — это идеализированный объект, который полностью поглощает
всё падающее на него электромагнитное излучение. Практической реализацией такого
объекта может служить полость (ящик) с малым отверстием\footnote{Размеры отверстия считаются малыми по
сравнению с размерами полости.}, стенки которой поддерживаются при постоянной
температуре.

Излучение, попавшее внутрь через отверстие, после многократных отражений
практически полностью поглощается стенками полости\footnote{Благодаря тому, что размеры
отверстия малы по сравнению с размерами полости, излучение не может выйти наружу
даже после нескольких переотражений от стенок.}. В состоянии теплового равновесия
устанавливается динамический баланс между процессами поглощения и испускания
излучения стенками полости. Часть этого равновесного излучения выходит наружу через
отверстие и может быть исследована экспериментально.

Для теоретического описания излучения абсолютно чёрного тела рассмотрим
ящик с идеально проводящими стенками и размерами $L_x$, $L_y$ и
$L_z$. В такой полости могут существовать только определённые моды
электромагнитных колебаний — стоячие волны с компонентами волнового вектора:

\begin{equation}
\begin{split}
|k_x| &= \frac{\pi n_x}{L_x}, \\
|k_y| &= \frac{\pi n_y}{L_y}, \\
|k_z| &= \frac{\pi n_z}{L_z},
\end{split}
\end{equation}

где $n_x$, $n_y$, $n_z$ — положительные целые числа, определяющие число полуволн,
укладывающихся вдоль соответствующего измерения.


Для каждой тройки чисел $(n_x, n_y, n_z)$ существуют две независимые поляризации
электромагнитной волны\footnote{Напомним, что электромагнитная волна может быть
  разложена на две поляризационные составляющие, например, с левой и правой
  круговой поляризацией.}. Таким образом, плотность энергии излучения в полости
можно представить как сумму по всем модам:

\[
\int \dd{W} = \frac{1}{V} \sum_{\xi=1}^{2} \sum_{\bm{n}} \bar{E}(\omega) = \frac{2}{V} \sum_{\bm{n}} \bar{E}(\omega),
\]

где $\bm{n} = (n_x, n_y, n_z)$, $V = L_x L_y L_z$ — объём полости, $\omega = kc = c\sqrt{k_x^2 + k_y^2 + k_z^2}$ — циклическая частота излучения, $\bar{E}(\omega)$ — средняя энергия, приходящаяся на одну моду колебаний.

Для макроскопической полости дискретный спектр частот становится практически
непрерывным\footnote{Если длины ребер ящика много больше характерных микроскопических
масштабов, то при изменении чисел $n_i$ на единицу модули компонент волновых
векторов $|k_i|$ меняются практически непрерывно.}. Это позволяет перейти от
суммирования по дискретным модам к интегрированию по волновым
векторам\footnote{При переходе от модулей компонет волновых векторов $|k_i|$ к
  компонентам волновых векторов $k_i$ в интеграле, мы учли, что $\omega$ не зависит
от знака этих компонет.}:

\[
\int \dd{W} = \frac{2}{V} \int_0^\infty \frac{L_x \dd{|k_x|}}{\pi} \int_0^\infty \frac{L_y \dd{|k_y|}}{\pi} \int_0^\infty \frac{L_z \dd{|k_z|}}{\pi} \bar{E}(\omega) = 2 \int \frac{\dd[3]{k}}{(2\pi)^3} \bar{E}(\omega)
\]

\[
\dd{W} = \frac{\bar{E}(\omega)}{4\pi^3} \dd[3]{k}
\]

% Предпоследний абзац
Теперь рассмотрим интенсивность излучения, выходящего через отверстие —
количество энергии, проходящей через единицу площади отверстия в единицу
времени. Выделим на поверхности полости бесконечно малый элемент площади $\Delta S$.
Излучение, которое проходит через этот элемент за время $\Delta t$, должно в
начальный момент времени находиться в непосредственной близости от него. Объём
пространства, из которого излучение со скоростью $c$ может достичь элемента
$\Delta S$ за время $\Delta t$, зависит от направления его распространения. Для излучения,
движущегося в направлении, составляющем угол $\theta$ с нормалью к элементу
$\Delta S$, этот объём представляет собой косой цилиндр с основанием $\Delta S$ и
образующей длиной $c\Delta t$ и равен $\Delta V = \Delta S \cdot c\Delta t \cdot \cos\theta$. Таким образом,
энергия $\dd{\mathcal{E}}$, переносимая излучением с волновыми векторами в интервале
$\dd[3]{k}$ через элемент $\Delta S$ за время $\Delta t$, равна произведению плотности
энергии $\dd{W}$ на этот объём:
\[
\dd{\mathcal{E}} = \dd{W} \cdot \Delta V = \frac{\bar{E}(\omega)}{4\pi^3} \dd[3]{k} \cdot \Delta S c \Delta t \cos\theta.
\]
Интенсивность, то есть энергия, проходящая через единицу площади в единицу времени, равна:
\[
\dd{I} = \frac{\dd{\mathcal{E}}}{\Delta S \Delta t} = \frac{\bar{E}(\omega)}{4\pi^3} c \cos\theta \dd[3]{k}.
\]

Для нахождения полной интенсивности излучения в интервале частот
$[\omega, \omega + \dd{\omega}]$, выходящего через отверстие, проинтегрируем $\dd{I}$ по всем
направлениям, из которых излучение может прийти к элементу площади. Для этого
перейдём в сферические координаты в $k$-пространстве. Учтём, что
$\dd[3]{k} = k^2 \dd{k} \dd{\Omega} = \frac{\omega^2}{c^3} \dd{\omega} \dd{\Omega}$, где
$\dd{\Omega} = \sin\theta \dd{\theta} \dd{\phi}$ — элемент телесного угла. Подставляя, получаем:
\[
\dd{I} = \frac{\bar{E}(\omega)}{4\pi^3} c \cos\theta \frac{\omega^2}{c^3} \dd{\omega} \dd{\Omega} = \frac{\bar{E}(\omega) \omega^2}{4\pi^3 c^2} \cos\theta \dd{\omega} \dd{\Omega}.
\]
Поскольку излучение распространяется изнутри полости наружу, вклад в
интенсивность дают только те лучи, для которых угол между направлением $\bm{k}$ и
нормалью к элементу площади $\Delta S$ удовлетворяет условию $0 \leqslant \theta < \pi/2$ (направление внутрь
полости соответствует $\theta > \pi/2$). Поэтому интегрирование по телесному углу
проводится по верхней полусфере ($\theta$ от $0$ до $\pi/2$, $\phi$ от $0$ до $2\pi$):
\[
\dv{I}{\omega} = \frac{\bar{E}(\omega) \omega^2}{4\pi^3 c^2} \int_0^{2\pi} \dd{\phi} \int_0^{\pi/2} \cos\theta \sin\theta \dd{\theta},
\]
откуда окончательно получаем искомое выражение для спектральной интенсивности:
\[
\dv{I}{\omega} = \frac{\bar{E}(\omega) \omega^2}{4\pi^2 c^2}.
\]
Аналогичное выражение для интенсивности, записанное через частоту $\nu$ (с учётом
соотношения $\omega = 2\pi\nu$), имеет вид:
\begin{equation}
  \label{eq:black-body-intensity}
  \dv{I}{\nu} = \frac{2\pi \bar{E}(2\pi\nu) \nu^2}{c^2}.
\end{equation}

\subsubsection*{Закон Рэлея-Джинса}

Согласно классической электродинамике энергия электромагнитного излучения принимает непрерывные значения. Для системы в тепловом равновесии при температуре $T$ распределение энергии по степеням свободы описывается законом Больцмана:

\[
\dv{\mathcal{W}}{E} = \frac{1}{Z} e^{-\beta E},
\]

где $\beta = 1/k_B T$, $k_B\approx 1.380 \cdot 10^{-16}\text{ эрг/К}$ — постоянная Больцмана, $Z = \int_0^\infty \dd{E} \exp(-\beta E)$ — статистическая сумма.

Средняя энергия, приходящаяся на одну моду колебаний:

\[
\bar{E}(\omega) = \frac{\int_0^\infty E e^{-\beta E} \dd{E}}{\int_0^\infty e^{-\beta E} \dd{E}} = -\pdv{\beta} \ln Z = \frac{1}{\beta} = k_B T
\]

Подстановка этого выражения в формулу для интенсивности~\eqref{eq:black-body-intensity} даёт закон Рэлея-Джинса~\cite{Rayleigh1900,Jeans1905}:
\begin{equation}
  \label{eq:rayleigh-jeans-law}
  \dv{I}{\nu} = \frac{2\pi k_B T \nu^2}{c^2}
\end{equation}

Полная интенсивность излучения получается интегрированием спектральной интенсивности по всем частотам:
\[
I = \int_0^\infty \dv{I}{\nu} \dd{\nu} = \frac{2\pi k_B T}{c^2} \int_0^\infty \nu^2 \dd{\nu}
\]
и оказывается расходящейся на верхнем пределе. Физически это означает, что вклад в полную интенсивность неограниченно растёт с увеличением частоты, что, конечно, противоречит экспериментальным данным по изучению абсолютно чёрного тела~\cite{Beckmann1898,Rubens1899,RubensKurlbaum1900,LummerPringsheim1900}. Закон Рэлея-Джинса согласуется с экспериментом только в области низких частот, а при высоких частотах наблюдается существенное расхождение. Это противоречие получило название <<ультрафиолетовой катастрофы>>.

К началу XX века проблема <<ультрафиолетовой катастрофы>> была хорошо известна, но рассматривалась как временная трудность, которую удастся устранить в рамках классической физики.

\subsubsection*{Формула Планка}

В 1900 году Макс Планк предложил решение этой проблемы. Он предположил, что энергия электромагнитного излучения испускается и поглощается не непрерывно, а дискретными порциями:

\[
E_n = n h \nu = n \hslash \omega, \quad n = 0, 1, 2, \dots
\]

где $h$ — новая фундаментальная постоянная (постоянная Планка), а $\hslash = h/2\pi$ — приведённая постоянная Планка, которая будет чаще использоваться в дальнейшем изложении курса квантовой механики.

В этом случае статистическая сумма для одной моды вычисляется как сумма геометрической прогрессии:

\[
Z = \sum_{n=0}^\infty e^{-\beta n \hslash \omega} = \frac{1}{1 - e^{-\beta \hslash \omega}}
\]

Средняя энергия на моду:

\[
\bar{E}(\omega) = \frac{\sum_{n=0}^\infty n \hslash \omega e^{-\beta n \hslash \omega}}{\sum_{n=0}^\infty e^{-\beta n \hslash \omega}} =
-\pdv{\beta} \ln Z = \frac{h\nu}{e^{\frac{h\nu}{k_B T}} - 1}
\]

Подстановка в формулу для интенсивности~\eqref{eq:black-body-intensity} даёт формулу Планка~\cite{Planck1900,Planck1901}:
\begin{equation}
\dv{I}{\nu} = \frac{2\pi h \nu^3}{c^2} \frac{1}{e^{\frac{h \nu}{k_B T}} - 1}
\end{equation}

Эта формула хорошо согласуется с экспериментальными данными во всей
области спектра и полностью разрешает проблему ультрафиолетовой катастрофы.
Полезно рассмотреть предельные случаи формулы Планка. При
$h\nu \ll k_{\mathrm B}T$ имеем
$e^{h\nu/(k_{\mathrm B}T)} - 1 \approx h\nu/(k_{\mathrm B}T)$, и формула
Планка переходит в закон Рэлея--Джинса~\eqref{eq:rayleigh-jeans-law}. С
другой стороны, при $h\nu \gg k_{\mathrm B}T$ знаменатель можно приближённо
заменить на $e^{h\nu/(k_{\mathrm B}T)}$, что даёт
\begin{equation}
  \dv{I}{\nu} \approx \frac{2\pi h \nu^3}{c^2} e^{-h\nu/(k_{\mathrm B}T)}.
\end{equation}
Таким образом, спектральная интенсивность экспоненциально убывает с ростом
частоты, что и устраняет ультрафиолетовую расходимость. Сравнение
спектральных зависимостей, следующих из закона Рэлея--Джинса и формулы
Планка, приведено на рис.~\ref{fig:planck-vs-rayleigh-jeans}.

\begin{figure}[h!]
  \centering
  \begin{tikzpicture}
    \begin{axis}[
      width=0.85\textwidth,
      height=0.55\textwidth,
      xmin=0, xmax=8,
      ymin=0, ymax=3.5,
      xlabel={$h\nu$ (в единицах $k_BT$)},
      ylabel={$\dv{I}{\nu}$ (отн. ед.)},
      samples=400,
      domain=0.01:8,
      axis lines=left,
      grid=major,
      grid style={gray!20},
      legend style={
        draw=none,
        at={(0.5,-0.18)},
        anchor=north,
        legend columns=-1
      },
      tick label style={font=\small},
      label style={font=\small},
    ]
      \addplot[line width=1pt, dashed, color=red!80!black] {x^2};
      \addlegendentry{закон Рэлея--Джинса}

      \addplot[line width=1pt, solid, color=blue!80!black]
        {x^3/(exp(x)-1)};
      \addlegendentry{формула Планка}

      \draw[dashed, gray] (axis cs:1,0) -- (axis cs:1,3.5);
      \node[anchor=south west, font=\scriptsize, text=gray]
        at (axis cs:1,3.5) {$h\nu=k_BT$};
    \end{axis}
  \end{tikzpicture}
  \caption{Спектральная интенсивность $\dv*{I}{\nu}$ как
  функция $h\nu$. По горизонтальной оси значения $h\nu$ отложены в единицах
  $k_BT$. Штриховая линия --- закон Рэлея--Джинса, сплошная ---
  формула Планка. Максимум спектральной интенсивности, задаваемой формулой
  Планка, достигается при $h\nu \approx 2.82 k_BT \sim k_BT$. При
  $h\nu\ll k_BT$ обе кривые совпадают. При
  $h\nu\gg k_BT$ спектральная интенсивность, задаваемая формулой
  Планка, экспоненциально убывает с ростом частоты. Спектральная
  интенсивность, задаваемая законом Рэлея--Джинса, напротив, неограниченно
  возрастает.}
  \label{fig:planck-vs-rayleigh-jeans}
\end{figure}

Сопоставление экспериментальных данных по излучению абсолютно чёрного тела с формулой Планка
позволило определить значение постоянной Планка:
\[
h \approx 6.626 \times 10^{-27} \text{ эрг·с}, \quad \hslash \approx 1.055 \times 10^{-27} \text{ эрг·с}\approx 6.582 \times 10^{-16} \text{ эВ·с}.
\]

Важно понимать, что в конце XIX века господствовали представления о волновой природе света,
основанные на успехах максвелловской теории электромагнитного поля. Большинство физических
явлений, связанных со светом, успешно описывались в рамках волновой теории. Поэтому
предположение Планка о квантовании энергии электромагнитного излучения первоначально воспринималось
как математический приём, позволяющий устранить ультрафиолетовую расходимость.
Считалось, что со временем удастся найти классическое объяснение спектра излучения
без введения дискретности.

Однако последующие исследования, особенно работы Эйнштейна по фотоэффекту, показали,
что квантование энергии носит фундаментальный характер и отражает корпускулярные (частицеподобные)
свойства электромагнитного излучения.

\subsection{Фотоэффект}

Следующее подтверждение корпускулярных свойств света было получено при изучении
явления фотоэффекта. Фотоэффект --- это процесс испускания электронов из вещества
под действием электромагнитного излучения.

В ранних экспериментальных исследованиях фотоэффекта (Герц, 1887 г.~\cite{Hertz1887}; Ленард,
1900 г.~\cite{Lenard1900}) были обнаружены закономерности, которые не удавалось
объяснить в рамках классической электродинамики. Рассмотрим схему эксперимента,
позволившего установить эти закономерности.

На рис.~\ref{fig:photoelectric-apparatus} приведена схема эксперимента по
изучению фотоэффекта. Свет определённой частоты $\nu$ проходит через кварцевое
окно и падает на металлическую пластину $C$ (фотокатод), помещённую в
вакуумную трубку. В результате фотоэффекта из пластины $C$ вылетают электроны.
Между фотокатодом $C$ и анодом $B$ с помощью внешнего источника создаётся
регулируемое напряжение $V$. Сила фототока регистрируется амперметром $A$.
Для измерения максимальной кинетической энергии фотоэлектронов на фотокатод
подаётся положительный потенциал, а на анод --- отрицательный, то есть
создаётся запирающее напряжение.

\begin{figure}[h!]
  \centering
  \begin{tikzpicture}[scale=0.9, >=stealth]
    \def\r{0.3}

    % Vacuum tube walls
    \draw[line width=3pt]
      (\r, 0) -- (10-\r, 0)
      arc[start angle=270, end angle=360, radius=\r]
      -- (10, 3-\r)
      arc[start angle=0, end angle=90, radius=\r]
      -- (7.5, 3);

    \draw[line width=3pt]
      (5.5, 3) -- (\r, 3)
      arc[start angle=90, end angle=180, radius=\r]
      -- (0, \r)
      arc[start angle=180, end angle=270, radius=\r];

    % Quartz window
    \draw[blue!50!black, line width=1.5pt] (7.5, 3) -- (5.5, 3);

    % Window label
    \node[anchor=south, font=\small, text=blue!50!black] at (8.7, 3.15) {кварцевое окно};

    % Tube label
    \node[anchor=north, font=\small] at (5.0, -0.15) {Вакуумная трубка};

    % Photocathode C (right) - with "+"
    \draw[line width=3pt] (8, 0.5) -- (8, 2.5);
    \node[above, font=\small] at (8, 2.45) {$+$};
    \node[right, font=\small] at (8.35, 1.5) {$C$};

    % Anode B (left) - with "-"
    \draw[line width=3pt] (2, 0.5) -- (2, 2.5);
    \node[above, font=\small] at (2, 2.45) {$-$};
    \node[left, font=\small] at (1.65, 1.5) {$B$};

    % Light: three parallel wavy lines with snake decoration
    \foreach \yzero in {5.5, 6.0, 6.5} {
      \draw[decorate, decoration={snake, amplitude=0.7mm, segment length=3mm},
            red!70!black, thick]
        (3.5, \yzero) -- (7.9, \yzero - 4.4);
    }
    \node[anchor=east, font=\small, text=red!70!black] at (3.3, 6.2) {Свет частоты $\nu$};

    % Electrons
    \foreach \x/\y in {2.7/1.2, 3.4/2.1, 4.1/0.8, 4.7/1.7, 5.3/2.3,
                       5.9/1.0, 3.0/2.4, 4.4/2.3, 6.3/0.7} {
      \fill[blue!70!black] (\x, \y) circle (1.5pt);
      \draw[->, blue!70!black, thick] (\x-0.1, \y) -- (\x-0.6, \y);
      \node[above right, font=\tiny, text=blue!70!black, inner sep=0.5pt]
        at (\x, \y) {$e^-$};
    }

    % External circuit
    \draw[thick] (8, 0.5) -- (8, -2.5);
    \draw[thick] (2, 0.5) -- (2, -2.5);

    % Battery with variable arrow
    \draw[thick] (8, -2.5) -- (6.5, -2.5);
    \draw[thick] (6.5, -2.1) -- (6.5, -2.9);
    \draw[thick] (6.3, -2.2) -- (6.3, -2.8);
    \draw[thick] (6.3, -2.5) -- (5.0, -2.5);
    \draw[->, thick] (6.0, -2.9) -- (6.8, -2.1);

    % Ammeter
    \draw[thick, fill=white] (4.5, -2.5) circle (0.45);
    \node[font=\small] at (4.5, -2.5) {$A$};

    \draw[thick] (4.05, -2.5) -- (2, -2.5);

    % Voltage label
    \draw[<->, thick] (2, -1.5) -- (8, -1.5);
    \node[below, font=\small] at (5.3, -1.55) {$V$};
  \end{tikzpicture}
  \caption{Схема экспериментальной установки для изучения фотоэффекта.
  Свет частоты $\nu$ проходит через кварцевое окно и освещает металлическую
  пластинку фотокатода $C$; вылетающие электроны
  собираются на аноде $B$. Фотокатод и анод соединены с источником регулируемого
  напряжения $V$. Сила фототока измеряется амперметром $A$.}
  \label{fig:photoelectric-apparatus}
\end{figure}

Первая закономерность фотоэффекта, которую не удавалось объяснить, заключалась в
следующем. Согласно классической электродинамике, свет представляет собой
электромагнитные волны. В рамках классической теории электрон должен был бы
накапливать энергию излучения постепенно так, что для накопления
необходимой для выхода из металла энергии ему понадобилось бы некоторое
заметное время. Однако эксперимент показывает, что фотоэлектроны появляются
практически мгновенно.

При увеличении запирающего напряжения электрическое поле сильнее тормозит
вылетевшие электроны. При некотором значении запирающего напряжения $V_0$
фототок обращается в нуль. Это означает, что даже самые быстрые фотоэлектроны
не могут преодолеть тормозящую разность потенциалов, так что максимальная
кинетическая энергия фотоэлектронов связана с величиной $V_0$ соотношением
\begin{equation}
  e V_0 = K_{\text{max}}.
\end{equation}
Вторая закономерность фотоэффекта связана с зависимостью величины $V_0$ от
свойств падающего света. Измерения показывают, что величина $V_0$ не зависит
от интенсивности света, но линейно зависит от его частоты
(рис.~\ref{fig:stopping-potential-frequency}). Следовательно, максимальная
кинетическая энергия фотоэлектронов \emph{линейно растёт с увеличением
частоты} излучения и \emph{не зависит от его интенсивности}. Наклон этой
линейной зависимости не зависит от материала катода.

\begin{figure}[h!]
  \centering
  \begin{tikzpicture}
    \begin{axis}[
      width=0.8\textwidth,
      height=0.55\textwidth,
      xmin=0, xmax=10,
      ymin=0, ymax=2.5,
      xlabel={Частота $\nu$ ($10^{14}$ Гц)},
      ylabel={Запирающее напряжение $V_0$ (В)},
      axis lines=left,
      axis line style={->},
      grid=major,
      grid style={gray!20},
      samples=100,
      domain=4.6:10,
      clip=false,
    ]
      \addplot[thick, blue!70!black] {0.4136*x - 1.9};

      % Threshold frequency for cesium
      \draw[dashed, gray] (4.6, 0) -- (4.6, 0.4);
      \node[below, font=\small] at (4.6, -0.1) {$\nu_0$};

      % Red boundary label (moved further left and down)
      \node[anchor=west, font=\small, text=gray] at (2.15, 0.52) {Красная граница};

      % Slope annotation (moved slightly left)
      \node[font=\small, text=blue!70!black] at (7.1, 1.4) {$\tg\alpha = h/e$};
    \end{axis}
  \end{tikzpicture}
  \caption{Зависимость запирающего напряжения $V_0$ от частоты падающего
  света $\nu$ для цезия. Тангенс угла наклона прямой равен отношению
  постоянной Планка к заряду электрона, $h/e$. Точка пересечения прямой с
  осью абсцисс определяет красную границу фотоэффекта $\nu_0 \approx
  4{,}6 \cdot 10^{14}$~Гц (что соответствует длине волны около 650~нм).}
  \label{fig:stopping-potential-frequency}
\end{figure}

Третья закономерность состоит в том, что существует такая частота $\nu_0$,
ниже которой фотоэффект не наблюдается. На рис.~\ref{fig:stopping-potential-frequency} эта частота соответствует
такой частоте излучения, при которой величина максимальной кинетической энергии
фотоэлектронов $K_{\text{max}} = eV_0$ обращается в нуль. Эта частота называется \emph{красной
границей фотоэффекта}. Величина этой частоты зависит только от материала фотокатода.

Таким образом, экспериментально установлены следующие основные свойства
фотоэффекта:
\begin{itemize}
\item \emph{мгновенность}: между облучением и вылетом электронов нет заметной
  задержки;
\item \emph{линейная зависимость максимальной кинетической энергии от
  частоты}: $K_{\text{max}}$ растёт с увеличением $\nu$;
\item \emph{независимость энергии фотоэлектронов от интенсивности света}:
  интенсивность влияет только на число вылетающих электронов;
\item \emph{существование красной границы}: при $\nu < \nu_0$ фотоэффект
  отсутствует.
\end{itemize}

В 1905 году Альберт Эйнштейн предположил, что квантование
энергии света носит фундаментальный характер и связано с тем, что свет при определённых условиях может рассматриваться как поток квантов
электромагнитного излучения (фотонов). Опираясь на эту гипотезу,
Эйнштейн опубликовал работу~\cite{Einstein1905}, в которой теоретически описал явление фотоэффекта.
Согласно этой работе, энергия кванта электромагнитного излучения
пропорциональна его частоте, $E = h\nu$, а процесс фотоэффекта заключается во
взаимодействии одиночного фотона с атомным электроном:
$\gamma + \text{атом} \to e^- + \text{ион}^+$.
При таком акте поглощения энергия фотона $h\nu$ передаётся электрону. Часть этой
энергии, обозначаемая $w$, тратится на преодоление сил, удерживающих данный
электрон внутри вещества. Эта работа $w$ может различаться для разных электронов
в зависимости от их начального энергетического состояния и потерь энергии на
пути к поверхности. Минимальное значение работы, необходимой для выхода
электрона из вещества, называется \emph{работой выхода} $w_0$ и является
характеристикой материала. Максимальную кинетическую энергию $K_{\text{max}}$
имеют электроны, для которых $w = w_0$. Для таких электронов справедливо
соотношение:
\[
K_{\text{max}} = h\nu - w_0.
\]
Из этого уравнения непосредственно следуют все экспериментально наблюдаемые особенности фотоэффекта:
\begin{itemize}
\item красная граница возникает при условии $K_{\text{max}} = 0$, что соответствует $h\nu_0 = w_0$;
\item линейная зависимость кинетической энергии от частоты: $K_{\text{max}} = h\nu - w_0$;
\item независимость энергии электронов от интенсивности света;
\item мгновенный характер процесса, так как энергия передаётся дискретными порциями.
\end{itemize}

Теория фотоэффекта, предложенная Эйнштейном, была экспериментально проверена
Милликеном в 1914 г.~\cite{Millikan1914}. Константа пропорциональности $h$ между частотой и энергией
кванта электромагнитного излучения совпала с постоянной Планка, определённой из
спектра излучения абсолютно чёрного тела.

Таким образом, фотоэффект стал одним из ключевых экспериментальных доказательств
квантовой природы света. Он наглядно продемонстрировал, что энергия
электромагнитного излучения поглощается дискретными порциями (квантами), и что
именно частота света, а не его интенсивность, определяет энергию этих порций.
Однако само по себе явление фотоэффекта не доказывает напрямую существование
фотонов как частиц, обладающих не только энергией, но и импульсом. Прямым
доказательством корпускулярных свойств электромагнитного излучения, то есть
существования фотонов, является \emph{эффект Комптона} (рассеяние рентгеновских
лучей на свободных электронах), в котором наблюдается упругое столкновение фотона
с электроном, подчиняющееся законам сохранения энергии и импульса.

\subsection{Эффект Комптона}

Эффект Комптона был открыт в 1923 году~\cite{Compton1923}.
В своем эксперименте Артур Комптон наблюдал рассеяние пучка рентгеновского
излучения на графитовой мишени.
Энергия фотона рентгеновского излучения много больше энергии связи электрона с
атомами в кристаллической решетке мишени. Следовательно, электроны, участвующие
в процессе рассеяния, могут рассматриваться как свободные. Процесс рассеяния
фотона на свободном электроне называется \emph{эффектом Комптона}.

С точки зрения классической электродинамики, картина рассеяния выглядит
следующим образом. Под действием электрического поля падающей электромагнитной
волны частоты $\nu_0$ электрон приходит в колебательное движение с той же частотой.
Колеблющийся электрон, в свою очередь, излучает электромагнитные волны, и
частота этого излучения равна частоте вынуждающей силы, то есть
$\nu_0$. Таким образом, классическая теория предсказывает, что длина волны
излучения $\lambda_0$ при рассеянии изменяться не должна.

Однако экспериментальные данные противоречат этому выводу. В спектре рассеянного
излучения наблюдается не один, а два пика (рис.~\ref{fig:compton-spectra}).
Первый пик соответствует излучению с неизменной длиной волны $\lambda = \lambda_0$
и связан с процессом \emph{релеевского рассеяния}.
Второй пик соответствует излучению с длиной волны $\lambda = \lambda_0 + \Delta\lambda$,
то есть наблюдается сдвиг длины волны $\Delta{\lambda}$. Экспериментальное
обнаружение этого пика привело к открытию \emph{эффекта Комптона}.
Величина смещения $\Delta{\lambda}$ не зависит от материала мишени и первоначальной
длины волны $\lambda_0$, а зависит исключительно от угла рассеяния $\theta$, то есть
угла между направлениями распространения падающего и рассеянного фотонов.

\begin{figure}[h!]
  \centering
  \begin{tikzpicture}
    \begin{axis}[
      width=0.65\textwidth,
      height=0.75\textwidth,
      xmin=0.69, xmax=0.77,
      ymin=0, ymax=7.5,
      axis lines=left,
      axis line style={->},
      xlabel={$\lambda$ (\AA)},
      ylabel={Интенсивность},
      xtick={0.70, 0.73, 0.75},
      extra x ticks={0.708},
      extra x tick labels={$\lambda_0$},
      ytick=\empty,
      samples=400,
      domain=0.69:0.77,
      clip=false,
    ]
      % Спектры при различных углах рассеяния (лоренцевские пики)
      \addplot[thick, blue!70!black, smooth]
        {1/(1 + ((x-0.708)/0.0035)^2)};

      \addplot[thick, blue!70!black, smooth]
        {1/(1 + ((x-0.708)/0.0035)^2)
         + 0.5/(1 + ((x-0.7151)/0.0035)^2) + 1.9};

      \addplot[thick, blue!70!black, smooth]
        {1/(1 + ((x-0.708)/0.0035)^2)
         + 0.5/(1 + ((x-0.7323)/0.0035)^2) + 3.8};

      \addplot[thick, blue!70!black, smooth]
        {1/(1 + ((x-0.708)/0.0035)^2)
         + 0.5/(1 + ((x-0.7495)/0.0035)^2) + 5.7};

      % Единая вертикальная штриховая линия, отмечающая положение lambda_0
      \draw[dashed, gray] (axis cs:0.708, 0) -- (axis cs:0.708, 7.5);

      % Подписи углов рассеяния
      \node[anchor=west, font=\small] at (axis cs:0.73, 0.5)
        {$\theta = 0^\circ$};
      \node[anchor=west, font=\small] at (axis cs:0.73, 2.4)
        {$\theta = 45^\circ$};
      \node[anchor=west, font=\small] at (axis cs:0.73, 4.75)
        {$\theta = 90^\circ$};
      \node[anchor=west, font=\small] at (axis cs:0.73, 6.2)
        {$\theta = 135^\circ$};
    \end{axis}
  \end{tikzpicture}
  \caption{Спектры рассеянного рентгеновского излучения при различных углах
  рассеяния $\theta$. По горизонтальной оси отложена длина волны $\lambda$
  в ангстремах, по вертикальной --- интенсивность рассеянного излучения.
  При $\theta = 0^\circ$ наблюдается только несмещённый пик (соответствует релеевскому рассеянию).
  По мере увеличения угла рассеяния появляется второй пик (соответствует эффекту
  Комптона), смещённый в сторону больших длин волн.}
  \label{fig:compton-spectra}
\end{figure}

Наличие этого сдвига может быть объяснено только в рамках корпускулярной теории.
Рассмотрим столкновение фотона с покоящимся свободным электроном. Кинетической
энергией электрона до столкновения пренебрегаем, так как она порядка энергии
связи, которая много меньше энергии рентгеновского фотона. Запишем законы
сохранения энергии и импульса для этого процесса. Закон сохранения энергии:
\begin{equation}
\label{eq:compton-energy}
p_0 c + m_e c^2 = p_1 c + E_e,
\end{equation}
где $p_0$ --- импульс налетающего фотона, $p_1$ --- импульс рассеянного фотона,
$m_e$ --- масса покоя электрона, $E_e$ --- полная энергия электрона отдачи. Закон
сохранения импульса:
\begin{equation}
\label{eq:compton-momentum}
\bm{p_0} = \bm{p_1} + \bm{p_e},
\end{equation}
где $\bm{p_0}$ и $\bm{p_1}$ --- векторы импульсов падающего и рассеянного фотонов,
$\bm{p_e}$ --- вектор импульса электрона после столкновения. Полная энергия
электрона $E_e$ связана с его импульсом $\bm{p_e}$ релятивистским соотношением:
\begin{equation}
\label{eq:electron-energy}
E_e^2 = p_e^2 c^2 + m^2_e c^4.
\end{equation}
Выразим $E_e$ из уравнения~\eqref{eq:compton-energy} и подставим
в~\eqref{eq:electron-energy}. После раскрытия квадрата и алгебраических
преобразований получим:
\begin{equation}
p_0^2 + p_1^2 - 2 p_0 p_1 + 2 m_e c (p_0 - p_1) = p_e^2.
\end{equation}
С другой стороны, из закона сохранения импульса~\eqref{eq:compton-momentum} следует:
\begin{equation}
p_e^2 = (\bm{p_0} - \bm{p_1})^2 = p_0^2 + p_1^2 - 2 p_0 p_1 \cos\theta,
\end{equation}
где $\theta$ --- угол между векторами $\bm{p_0}$ и $\bm{p_1}$.
Приравнивая оба выражения для $p_e^2$, после сокращения одинаковых членов получаем:
\begin{equation}
- 2 p_0 p_1 + 2 m_e c (p_0 - p_1) = - 2 p_0 p_1 \cos\theta.
\end{equation}
Преобразуя это выражение, приходим к уравнению:
\begin{equation}
\frac{1}{p_1} - \frac{1}{p_0} = \frac{1}{m_e c} (1 - \cos\theta).
\end{equation}
Учитывая связь импульса фотона с длиной волны $p = h/\lambda$, получаем окончательную
формулу для комптоновского сдвига:
\begin{equation}
  \label{eq:compton-shift}
  \Delta\lambda = \lambda_1 - \lambda_0 = \frac{h}{m_e c} (1 - \cos\theta) = 2\lambda_C \sin^2\frac{\theta}{2},
\end{equation}
где $\lambda_C = h/(m_e c)$ --- комптоновская длина волны электрона.

Наличие несмещенной линии (релеевского рассеяния) объясняется рассеянием фотонов на электронах, сильно связанных в атомах. В этом случае фотон рассеивается на атоме как целом. Для описания такого процесса в формуле для комптоновского сдвига~\eqref{eq:compton-shift}
массу электрона $m_e$ следует заменить на массу атома $M$. Поскольку масса атома
$M$ много больше массы электрона $m_e$, величина сдвига $\Delta\lambda$ становится ничтожно
малой и практически не наблюдается.

\subsection{Опыт Резерфорда}
В 1911 г. Эрнест Резерфорд провёл эксперимент по рассеянию
$\alpha$-частиц\footnote{$\alpha$-частица --- ядро атома гелия-4, состоящее из двух протонов и
  двух нейтронов.} на мишени из золотой фольги~\cite{Rutherford1911}. В эксперименте наблюдалось, что
хотя большинство $\alpha$-частиц проходило через фольгу с незначительным отклонением,
некоторые частицы рассеивались на большие углы $\sim 1 \text{ рад}$.
Такое рассеяние возможно только если весь положительный заряд и почти вся масса
атома сосредоточены в малой области (атомном ядре) в центре атома. Для
объяснения рассеяния $\alpha$-частиц на углы $\sim 1\text{ рад}$ радиус атомного ядра
должен составлять $\sim 10^{-12}\text{ см}$ по порядку величины.
На основании этих результатов Резерфорд предложил планетарную (ядерную) модель
атома: тяжёлое положительно
заряженное ядро, вокруг которого движутся электроны. В рамках этой модели
характерный размер атома определяется радиусом орбиты внешних электронов
и составляет $\sim 10^{-8}\text{ см}$, что значительно превышает размер ядра.
Однако планетарная модель в рамках классической физики приводит к противоречию:
электроны на орбитах движутся с ускорением и должны излучать электромагнитные
волны, теряя энергию. В результате за время $\sim 10^{-11}\text{ с}$ электроны
должны упасть на ядро, что означало бы нестабильность атома. Устойчивость
атомов, наблюдаемая в повседневном опыте, указывает на неприменимость
классических законов к микроскопическим объектам.

\subsection{Атомные спектры}
Классические законы физики, применимые к макроскопическим телам, оказываются
неспособны объяснить наблюдаемые атомные спектры. Согласно классической
электродинамике, электрон, движущийся по орбите вокруг атомного ядра, должен
излучать непрерывный спектр. Однако экспериментально наблюдается дискретный
линейчатый спектр --- набор отдельных спектральных линий.

Эмпирические закономерности в спектре атома водорода были установлены Бальмером
(1885 г.)~\cite{Balmer1885} и Ридбергом (1890 г.)~\cite{Rydberg1890}. Положение спектральных линий описывается
формулой:
\begin{equation}
\frac{1}{\lambda} = R_H\left(\frac{1}{m^2} - \frac{1}{n^2}\right),
\end{equation}
где $m$ и $n$ --- натуральные числа ($n > m$),
$R_H \approx 10967757.6\pm1.2\text{ м}^{-1}$ --- постоянная Ридберга для водорода. Для
перехода к частоте фотона формулу следует умножить на скорость света $c$.
Энергия излученного фотона выражается как:
\begin{equation}
E = \hslash\omega = \text{Ry}\left(\frac{1}{m^2} - \frac{1}{n^2}\right),
\end{equation}
где постоянная Ридберга в энергетических единицах $\text{Ry} \approx 13.6\text{ эВ}$.

Совокупность спектральных линий при фиксированном $m$ образует спектральную
серию. Основные серии атома водорода:
\begin{itemize}
\item серия Лаймана ($m=1$) --- ультрафиолетовая область~\cite{Lyman1906};
\item серия Бальмера ($m=2$) --- видимая и ближняя ультрафиолетовая области~\cite{Balmer1885};
\item серия Пашена ($m=3$) --- инфракрасная область~\cite{Paschen1908};
\item серии Брэккета ($m=4$)~\cite{Brackett1922} и Пфунда ($m=5$)~\cite{Pfund1924} --- дальняя инфракрасная область.
\end{itemize}

\subsection{Модель атома Бора}
Для решения проблем стабильности атома и объяснения дискретных атомных спектров
Нильс Бор в 1913 году~\cite{Bohr1913} предложил модель атома, основанную на специальных
постулатах. Эта модель была разработана в рамках так называемой <<старой
квантовой теории>> — подхода к описанию квантовых явлений, развивавшегося до
создания современной квантовой механики. Модель Бора основывается на следующих
постулатах:
\begin{enumerate}
\item Электроны движутся по стационарным круговым орбитам вокруг ядра под действием
  кулоновского притяжения, подчиняясь законам классической механики.
\item Допустимы не все орбиты, а только те, для которых орбитальный момент
  количества движения электрона квантуется:
\begin{equation}
L = m_evr = n\hslash,
\end{equation}
где $n$ --- целое положительное число ($n = 1, 2, 3, \ldots$).
\item Находясь на стационарной орбите, электрон не излучает, несмотря на наличие ускорения.
\item При переходе электрона между стационарными орбитами происходит излучение или
  поглощение фотона. Частота фотона определяется разностью энергий орбит:
\begin{equation}
\hslash\omega = E_i - E_f.
\end{equation}
\end{enumerate}

Используя эти постулаты, выведем выражение для уровней энергии водородоподобного
иона\footnote{Водородоподобный ион --- это система, состоящая из атомного ядра с
  положительным зарядом $Ze$, вокруг которого движется единственный электрон.}.
Рассмотрим электрон на круговой орбите вокруг ядра с зарядом $Ze$. Условие
равновесия между кулоновской и центробежной силами:
\begin{equation}
\frac{Ze^2}{r^2} = \frac{m_ev^2}{r}.
\end{equation}
Из условия квантования момента импульса $m_evr = n\hslash$ находим радиус разрешенных орбит:
\begin{equation}
r_n = \frac{\hslash^2n^2}{m_eZe^2}.
\end{equation}
Полная энергия электрона на $n$-й орбите:
\begin{equation}
E_n = \frac{m_ev^2}{2} - \frac{Ze^2}{r} = -\frac{m_eZ^2e^4}{2\hslash^2n^2}.
\end{equation}
Интересно, что это выражение, полученное в рамках модели Бора, в точности
совпадает с результатом, получаемым в современной квантовой механике.
Для атома водорода ($Z=1$) основное состояние соответствует $n=1$ с энергией
$E_1 = -13.6$ эВ. Постоянная Ридберга выражается через параметры теории:
\begin{equation}
\text{Ry} = \frac{m_ee^4}{2\hslash^2} \approx 13.6\text{ эВ}.
\end{equation}

Несмотря на успех в объяснении спектра водорода, модель Бора имела фундаментальные ограничения:
\begin{itemize}
\item Неприменима к описанию многоэлектронных атомов
\item Не объясняла интенсивности спектральных линий
\item Не давала понимания механизма переходов между уровнями
\item Содержала внутреннее противоречие: сочетание классических орбит с квантовыми постулатами
\end{itemize}

Эти ограничения привели к необходимости разработки современной квантовой
механики, основанной на принципиально новых концепциях.

\section{Гипотеза Луи де Бройля}

В 1924 г. Луи де Бройлем была высказана гипотеза о волновых свойствах частиц
вещества~\cite{deBroglie1924, deBroglie1925}. Он утверждал, что каждой частице с энергией $E$ и импульсом $\bm{p}$
соответствует волна с частотой $\omega$ и волновым вектором $\bm{k}$ такими, что
$\omega=E/\hslash$, $\bm{k} = \bm{p}/\hslash$. Наличие у частиц волновых свойств было
подтверждено многочисленными опытами по дифракции и интерференции пучков
различных частиц, например, таких как электроны и нейтроны. Таким образом,
частицы материи могут проявлять как корпускулярные свойства, так и волновые. Это
обстоятельство называется \textit{корпускулярно-волновым дуализмом}.

Для иллюстрации условий, при которых частицы проявляют корпускулярные или
волновые свойства, рассмотрим свободный электрон с энергией $E=10\text{ эВ}$.
Поскольку электрон свободный, то его энергия совпадает с кинетической энергией,
то есть $E=p^2/2m_e$. Из последнего соотношения можно выразить модуль импульса
электрона, $p = \sqrt{2m_eE}$. Далее запишем связь длины волны с импульсом
\begin{equation}
  \label{eq:wave-length}
  \begin{split}
    \lambda &= \frac{2\pi}{k} = \frac{2\pi\hslash}{p} = \frac{2\pi\hslash}{\sqrt{2m_eE}} = \frac{2\pi\hslash c}{\sqrt{2m_ec^2E}}  \approx \\
      &\approx\frac{2\cdot{3.14}\cdot{6.58}\times{10^{-16}}\text{ эВ}\cdot\text{сек}\cdot{3}\times{10^{10}}\frac{\text{см}}{\text{сек}}}{\sqrt{2 \cdot 0.5\times{10^6}\text{ эВ}\cdot{10}\text{ эВ}}}\sim 10^{-8}\text{ см}.
  \end{split}
\end{equation}

Из курса электродинамики известно, что на координатном масштабе, много большем
длины волны, для описания излучения применима геометрическая оптика. Если
масштаб, на котором рассматривается излучение, сопоставим с длиной волны, то
необходимо использовать волновое описание. Применяя эту аналогию к электрону,
делаем вывод, что мы можем считать, что электрон с энергией $10\text{ эВ}$
проявляет корпускулярные свойства, если мы рассматриваем его движение на
масштабе много большем, чем $\lambda=10^{-8}\text{ см}$. В противном случае необходимо
учитывать волновые свойства электрона.

Если рассматривать движение электрона с энергией $10\text{ эВ}$ на масштабах,
сопоставимых с размером атома, то следует отметить, что размер атома по порядку
величины равен $10^{-8}\text{ см}$, то есть того же порядка, что и длина волны
такого электрона. Поэтому при изучении движения электрона на атомных масштабах
мы будем наблюдать проявление его волновых свойств. Для протона той же энергии
($E=10\text{ эВ}$) длина волны окажется меньше, чем у электрона, за счет
квадратного корня из массы в знаменателе формулы~\eqref{eq:wave-length}. Это
означает, что волновые свойства протона будут наблюдаться на гораздо меньших
масштабах, чем у электрона. Для тяжелых атомов, молекул и других массивных
частиц волновые свойства будут проявляться на ещё меньших масштабах. Таким образом, на
практике волновых свойств у макроскопических объетов мы наблюдать не будем.

\printbibliography

\end{document}
